% Comparativa de las tres formas de resolver una pregunta de seguimiento
% (IT-106). Las cifras las produce scripts/experimento_conversacion.py sobre las
% 39 conversaciones derivadas de data/grados.json; si se regenera el corpus hay
% que volver a ejecutarlo y traer la salida real, no ajustar los numeros aqui.
\begin{table}[htbp]
  \centering
  \caption{\textcolor{borrador2}{Comparación de las tres formas de convertir una
      pregunta de seguimiento en una consulta buscable, sobre 39 conversaciones
      derivadas del propio corpus.}}
  \label{tab:conversacion}
  \small
  \begin{tabular}{@{}l cccc cccc@{}}
    \toprule
                                    & \multicolumn{4}{c}{\textbf{Unidad recuperada}} & \multicolumn{4}{c}{\textbf{MRR}}                                                                                                     \\
    \cmidrule(lr){2-5} \cmidrule(lr){6-9}
    \textbf{Estrategia}             & CC                                             & CT                               & SR    & \textbf{Total} & CC    & CT    & SR    & \textbf{Total} \\
    \midrule
    Sin conversación                & 0,080                                          & 0,333                            & 1,000 & 0,359          & 0,005 & 0,022 & 0,275 & 0,082          \\
    Concatenación de la anterior    & \textbf{1,000}                                 & 0,333                            & 1,000 & 0,949          & 0,653 & 0,022 & 0,275 & 0,498          \\
    \textbf{Sujeto deducido y filtrado} & \textbf{1,000}                             & \textbf{1,000}                   & \textbf{1,000} & \textbf{1,000} & \textbf{0,807} & \textbf{1,000} & \textbf{1,000} & \textbf{0,877} \\
    \bottomrule
  \end{tabular}
  \begin{minipage}{0.95\textwidth}
    \vspace{4pt}
    \tiny
    \textbf{Notas:} \textbf{CC}: la pregunta cambia el curso y mantiene la
    titulación (25 conversaciones). \textbf{CT}: cambia la titulación y mantiene
    lo que se pregunta (3). \textbf{SR}: la titulación la introdujo el asistente
    en su respuesta y no el estudiante en su pregunta (11).
    \textbf{Unidad recuperada}: proporción de conversaciones en las que el
    fragmento que responde a la pregunta entra entre los recuperados, con
    $K = 20$. \textbf{MRR}: media del inverso del puesto que ocupa ese fragmento;
    vale 1 si sale siempre el primero. \textcolor{borrador2}{Cifras producidas
      por \texttt{scripts/experimento\_conversacion.py} sobre el corpus de la
      extracción del 16/08/2026.}
  \end{minipage}
\end{table}
